\documentclass[11pt]{article}
\usepackage[a4paper,margin=1in]{geometry}
\usepackage{amsmath,amssymb,amsthm}
\usepackage{booktabs}
\usepackage{array}
\usepackage{hyperref}

\newtheorem{definition}{Definition}
\newtheorem{theorem}{Theorem}
\newtheorem{nogo}{No-Go}

\newcommand{\A}{\mathbb A}
\newcommand{\Q}{\mathbb Q}
\newcommand{\HC}{H_{\mathrm C}}
\newcommand{\DC}{D_{\mathrm C}}
\newcommand{\XiConnes}{\Xi_{\mathrm{Connes}}}
\newcommand{\CRCFT}{\mathrm{CRCFT}}

\title{Step 184: Connes Adelic / Noncommutative-Geometric Carrier Pivot}
\author{RH Membrane Adequacy Track}
\date{}

\begin{document}
\maketitle

\section*{Verdict}

\[
\boxed{\mathrm{connes\_verdict}=V_{\mathrm{connes\_CRCFT\_TE}}.}
\]

The Connes adelic/noncommutative-geometric carrier lands in
\(\CRCFT\)-TE.  The native full-strength trace-formula closure is the
classical Connes RH-equivalent spectral-realization condition.  No inherited
record supplies an independent completed spectral ledger, determinant
convergence theorem, or tail/exhaustivity record.

\section*{Inherited Record Audit}

\paragraph{Step 90.}
Trace-formula and automorphic/representation sources are recorded as
potentially genuine only if the chosen completed carrier is automorphic or
adelic and if Plancherel or trace-formula data provide a full anti-invariant
lower frame.  Otherwise they are public trace shadows.

\paragraph{Step 91.}
The character-source record states that a conductor ladder is a moving
finite-window source unless embedded into a fixed completed response space,
such as an adelic/idele-class character space, with a Plancherel or
tail-exhaustivity theorem.

\paragraph{Step 93.}
The Tier-0 triage imports Connes--Consani 2020 as the archimedean positivity
slot and identifies Connes--Consani--Moscovici semilocal spaces as the natural
arena for the cross-term.  It also records that finite spectral triples give
evidence but require convergence to the completed zero ledger.

\paragraph{Steps 99--100.}
The Plancherel/exhaustivity records state that finite windows and finite
source frames promote only with a fixed exhaustive tail form.  Step 100
instantiates a semilocal carrier
\[
Y_S^- = L^2(\mathbb R,dm_S)^-,
\qquad
dm_S(\xi)=\left|\prod_{v\in S}L_v(1/2-i\xi)\right|^2d\xi,
\]
but leaves the active analytic target as the completed source/tail domination
problem.  It does not construct the full Connes adelic spectral ledger.

\section*{Connes Carrier Declaration}

\begin{definition}[Typed Connes carrier]
Let
\[
X_\Q=\A_\Q/\Q^*,\qquad C_\Q=\A_\Q^*/\Q^*
\]
be the adele class space and idele class group.  The typed Connes carrier
\(\HC\) is the completed adelic/cohomological Hilbert carrier associated to
\(X_\Q\), with the trivial modes removed, on which \(C_\Q\) acts by scaling.
\end{definition}

The operator \(\DC\) is declared as the infinitesimal generator of the modulus
subgroup of the scaling action, or equivalently the generator whose
distributional trace is used on the spectral side of Connes' trace formula.
This declaration is typed rather than computational: the exact completed
domain, cokernel model, and regularization are classical Connes/Meyer records
and are not supplied by the inherited artifacts.

\section*{Trace Formula and Parent Residual}

For a suitable test function \(h\) on \(C_\Q\), the Connes trace formula has
the schematic shape
\[
\operatorname{Tr}\bigl(\pi_{\mathrm C}(h)\mid \HC\bigr)
=
\mathcal E_{\mathrm{arith}}(h),
\]
where the right side is the explicit-formula distribution: prime or
von-Mangoldt orbital terms, archimedean terms, and pole/trivial-zero
corrections.

Define
\[
\XiConnes(h)
:=
\operatorname{Tr}\bigl(\pi_{\mathrm C}(h)\mid \HC\bigr)
-
\mathcal E_{\mathrm{arith}}(h)
\]
after quotienting trivial modes and after imposing the required completed
tail/exhaustivity conditions.  Thus
\[
\XiConnes=0
\]
means that the completed Connes spectral ledger realizes the zeta zero ledger
with the full trace-formula positivity/effectivity required by the classical
Connes program.

\section*{CRCFT Mode Audit}

\begin{center}
\begin{tabular}{p{0.18\linewidth}p{0.18\linewidth}p{0.54\linewidth}}
\toprule
Audit & Result & Reason\\
\midrule
T3a CTMT & not primary
& No inherited carrier-native basis or kernel reduces \(\XiConnes\) to a
matrix-element family before the trace/ledger gate.\\
T3b TE & yes
& Full trace-formula/spectral-realization closure is the Connes RH-equivalent
condition at full strength.\\
T3c BF & secondary
& Bridges to Hecke, Burnol, or zeta residuals are missing, but the native
Connes route already lands in TE.\\
T3d non-CRCFT & none found
& Finite spectral triples and semilocal source routes require determinant or
zero-ledger convergence plus tail records.\\
\bottomrule
\end{tabular}
\end{center}

\section*{Initial Cascade}

\begin{center}
\begin{tabular}{p{0.18\linewidth}p{0.25\linewidth}p{0.34\linewidth}p{0.16\linewidth}}
\toprule
Branch & Name & Target & Status\\
\midrule
C-trace & completed trace-formula closure
& prove spectral trace equals arithmetic explicit-formula side over the full test class
& \(\CRCFT\)-TE\\
C-ledger & completed spectral ledger
& promote finite spectral triples or approximants to the completed zero ledger
& finite-window support only\\
C-bridge & bridge to prior carriers
& compare Connes to Hecke/Burnol/zeta residuals without scalar identity collapse
& bridge missing\\
C-source & idele character/source coercivity
& turn idele Plancherel sources into a fixed exhaustive lower frame
& tail required\\
\bottomrule
\end{tabular}
\end{center}

\begin{nogo}[Connes target-equivalence no-go]
Under inherited records, the native Connes adelic/NCG carrier exhibits
\(\CRCFT\)-TE: full closure of \(\XiConnes\) is target-equivalent to RH unless
an independent completed spectral ledger and tail/exhaustivity theorem are
supplied.
\end{nogo}

\begin{proof}
The carrier is genuinely native and adelic, so the primary failure is not
that the carrier is a scalar shadow.  However, the inherited records only
import semilocal and finite-window parts: archimedean positivity,
Hardy--Titchmarsh semilocal response spaces, finite spectral triples, and
Plancherel/tail promotion rules.  Each inherited finite or semilocal object is
explicitly scoped as support-only until a completed zero ledger and exhaustive
tail record are supplied.

The full Connes statement that the spectral trace on the completed adelic
carrier realizes the zeta zero ledger is exactly the classical
Connes/RH-equivalent closure condition.  Using that closure as a tool would
therefore be target-equivalence smuggling.  Since no inherited basis or kernel
gives a prior \(\CTMT\)-style matrix-element terminus, and no independent
carrier bridge supplies the missing ledger, the typed classification is
\(\CRCFT\)-TE.
\end{proof}

\section*{Classical Theorems Cited}

\begin{itemize}
\item A. Connes, \emph{Trace formula in noncommutative geometry and the zeros
of the Riemann zeta function}, Selecta Mathematica, 1999.
\item R. Meyer, work from 2005 onward on functional-analytic reformulations of
Connes' spectral realization.
\item C. Deninger's cohomological/motivic RH program, used here only as
context for the target-equivalence risk in cohomological spectral carriers.
\end{itemize}

\section*{Coverage Status}

\[
\boxed{\mathrm{CRCFT\_coverage\_conjecture\_status}=\mathrm{strengthens}.}
\]
The Connes carrier adds another \(\CRCFT\)-TE instance.  It does not refute
the CRCFT coverage conjecture because no non-\(\CRCFT\), non-target-equivalent
route is supplied by the inherited records.

\section*{Nonclaim Boundary}

This step does not prove RH, does not assume RH, and does not construct the
completed Connes spectral ledger.  It does not close \(\XiConnes\) or any
previous residual.  It records a target-equivalence typed no-go and preserves
all retained no-gos.

\end{document}
